\section{Exact entry-flip calculus and persistent discrete traps}
\label{sec:spectral-flip-traps}
Put $F(H)=\|HH^T-nI\|_F^2$ for an order-$n$ sign matrix.
This discrete energy landscape differs from optimization over the real
orthogonal group. Squared-Gram spin formulations and nonglobal local
minima are established background; see the primary article
\emph{Quantum computing formulation of some classical Hadamard matrix
searching methods and its implementation on a quantum computer},
\url{https://doi.org/10.1038/s41598-021-03586-0}.
The earlier primary article \emph{Finding a Hadamard Matrix by Simulated
Quantum Annealing} (2018), \url{https://www.mdpi.com/1099-4300/20/2/141},
also studies an energy-search formulation.
The explicit certificates and tensor bounds below are derived directly;
novelty is not asserted.

\begin{lemma}[Exact one- and two-entry changes]\label{lem:spectral-flip-calculus}
Let $G=HH^T$, $K=H^TH$. Flipping entry $(i,a)$ changes energy by
\[
 d_{ia}=8\big[(2n-1)-h_{ia}(GH)_{ia}\big].
\]
For two distinct flipped cells the change is the sum of their individual
changes plus the interaction
\[
 \begin{cases}
 16(h_{ia}h_{ib}K_{ab}-1),&\text{same row }i,\ a\ne b,\\
 16(h_{ia}h_{ja}G_{ij}-1),&\text{same column }a,\ i\ne j,\\
 16h_{ia}h_{ib}h_{ja}h_{jb},&i\ne j,\ a\ne b
 \quad\text{for cells }(i,a),(j,b).
 \end{cases}
\]
For even $n$, every single-entry energy change is eight times an odd
integer, so there are no neutral single-entry moves.
\end{lemma}
\begin{proof}
Only row-$i$ correlations change in a single flip:
$G_{ij}\mapsto G_{ij}-2h_{ia}h_{ja}$ for $j\ne i$.
Summing the changes of their squared values, twice for symmetry, gives
$8[(n-1)-h_{ia}\sum_{j\ne i}h_{ja}G_{ij}]$, the displayed formula.
Alternatively expand
\[
 F(H)=n^2(n-1)+4\sum_{i<j,\ a<b}h_{ia}h_{ib}h_{ja}h_{jb}.
\]
The two-flip interaction is sixteen times the sum of rectangle products
containing both cells. This gives the three cases. For even $n$ all
entries of $G$ are even, so $(GH)_{ia}$ is even while $2n-1$ is odd.
\end{proof}

\begin{theorem}[An exact small trap and its optimal escape barrier]
\label{thm:spectral-eight-trap}
Let $A$ be the $8\times8$ sign matrix whose negative-entry bitmasks in
successive rows are
\[
 (84,44,30,230,133,73,111,51),
\]
where bit $j$ denotes column $j$, starting at zero. Then $F(A)=152$,
every single-entry flip increases energy by at least eight, and every
two-entry flip has nonnegative energy change. A three-entry flip lowers
energy to $128$. The minimum Hamming distance from $A$ to any
strictly lower-energy matrix is exactly three. The least energy barrier
along single-entry paths to a lower-energy matrix is exactly eight.
\end{theorem}
\begin{proof}[Exact finite certificate]
The matrix of single-flip changes divided by eight is
\[
 \begin{pmatrix}
 1&5&5&9&5&1&9&9\\
 5&9&1&1&5&9&9&5\\
 1&9&5&5&1&9&5&1\\
 5&5&1&5&5&1&5&1\\
 9&1&9&5&5&5&9&1\\
 1&1&5&1&5&5&5&5\\
 9&5&9&1&1&1&5&5\\
 5&1&1&9&1&5&5&9
 \end{pmatrix}.
\]
The checker evaluates every one of the $\binom{64}{2}=2016$ two-cell
changes using both direct Gram energy and the interaction formula;
their minimum is zero. It also evaluates all $\binom{64}{3}=41664$
three-cell changes, whose minimum is $-24$. Explicitly, flipping cells
with flattened indices $0,5,60$ gives energies
$152\to160\to152\to128$. The first step of any single-entry path
costs at least eight, so this path attains the minimum possible barrier.
The one- and two-flip checks prove the claimed distance-three minimum.
All calculations are exact integer operations, not a heuristic search
certificate. The original randomized hill search merely found the
matrix; none of the assertions depends on its search history.
\end{proof}

\begin{theorem}[Infinite strict two-flip traps]
\label{thm:spectral-tensor-discrete-traps}
For any sign Hadamard matrix $B$ of order $m\ge2$, the matrix
$A\otimes B$ has positive energy $152m^3$ and strictly increases energy
under every change of one or two entries. The exact minimum increases are
\[
 \min\Delta_1=16m-8,\qquad \min\Delta_2=32(m-1).
\]
In particular this gives infinitely many explicit nonglobal strict
one- and two-entry local minima at known orders $8\cdot2^r$, $r\ge1$.
\end{theorem}
\begin{proof}
Let $Q=AA^T$. The lift has row Gram $Q\otimes mI_m$, and
$HH^TH=(Q A)\otimes mB$. If a base cell has single-flip change $d$,
its lifted cells have change $md+8(m-1)$. Since the minimum base change
is eight, the exact lifted minimum is $16m-8$.
Energy scales as
$\|(Q-8I_8)\otimes mI_m\|_F^2=m^3F(A)$.

For two cells sharing neither row nor column, the interaction is at
least $-16$, so their total increase is at least $32m-32$.
For two cells in the same lifted row, if their Hadamard column indices
differ, column-Gram orthogonality again makes the interaction $-16$.
If those indices agree, their base columns differ and the exact lifted
pair change is $m\Delta_{2,\mathrm{base}}+32(m-1)$, nonnegative in its
first term. The same argument applies to a shared lifted column.
Thus every pair costs at least $32(m-1)$. Equality is attained by lifting
the neutral base pair $(0,0),(0,5)$ with the same Hadamard row and column
indices. Positive energy proves nonglobality, since tensoring Walsh order eight with $B$ gives an order-$8m$
Hadamard matrix.
\end{proof}

For completeness, the checker exhausts all $512$ normalized order-four
sign matrices. Only the six Hadamard matrices are single-entry local
minima. Row and column sign normalization preserves energy and the
entry-flip graph, so order eight is the first admissible order at which
this greedy-descent obstruction occurs. The example does not rule out
nonlocal methods: its base has a certified three-flip escape. It does
rule out a proof that arbitrary one- or two-entry improving moves must
always reach a Hadamard matrix. No all-order construction is supplied.
